% BHU PYQ -- AIHC & Archaeology: Political History of South India (600 CE - 1300 CE) (2025-26 NEP)
\documentclass[11pt,a4paper]{article}
\usepackage[a4paper,top=2.5cm,bottom=2.5cm,left=2.8cm,right=2.8cm,headheight=14pt]{geometry}
\usepackage{amsmath,amssymb}
\usepackage{fontspec}
\setmainfont{Kohinoor Devanagari}
\usepackage{microtype,fancyhdr,enumitem,hyperref,lastpage,parskip}

\hypersetup{colorlinks=true,urlcolor=black,linkcolor=black,pdftitle={AIHMJ31: Political History of South India -- BHU NEP 2025-26}}

\pagestyle{fancy}\fancyhf{}
\renewcommand{\headrulewidth}{0.4pt}\renewcommand{\footrulewidth}{0.4pt}
\fancyhead[L]{\small \textit{Banaras Hindu University}}
\fancyhead[R]{\small \textit{AIHMJ31: Political History of South India (2025-26)}}
\fancyfoot[C]{\small Page \thepage\ of \pageref{LastPage}}

\newlist{parts}{enumerate}{2}
\setlist[parts,1]{label=(\alph*),leftmargin=*,itemsep=4pt,topsep=3pt}
\setlist[parts,2]{label=(\roman*),leftmargin=*,itemsep=2pt,topsep=2pt}

\newcommand{\pts}[1]{\hfill \textit{[#1]}}

\begin{document}

\begin{center}
  {\large\bfseries BANARAS HINDU UNIVERSITY}\\[2pt]
  {\normalsize B. A. (Hons.) Semester III Examination, 2025-26 (NEP)}\\[6pt]
  \rule{0.8\linewidth}{0.5pt}\\[6pt]
  {\large\bfseries ANCIENT INDIAN HISTORY, CULTURE \& ARCHAEOLOGY}\\[4pt]
  {\large\bfseries Paper Code: AIHMJ31 : Political History of South India (Circa 600 CE -- 1300 CE)}\\[4pt]
  {\normalsize Time: 2:30 Hours \hfill Full Marks: 70}\\[2pt]
  {\small REF NO. 2025-26/D-III/068}
\end{center}
\rule{\linewidth}{0.4pt}
\smallskip

\noindent \textit{Instructions: Attempt all questions. Questions of Section A, B and C carry 10, 05 and 02 marks each respectively. Write your Roll No.\ at the top immediately on receipt of this question paper.}

\smallskip
\noindent \textit{सभी प्रश्नों के उत्तर दीजिए। खण्ड अ, ब तथा स के प्रश्न क्रमशः 10, 05 एवं 02 अंकों के हैं।}

\bigskip

\section*{SECTION -- A / खण्ड -- अ}
\noindent \textit{Note: Long Answer Type Questions. Answer approximately in 250 words.}\pts{3 $\times$ 10 = 30}

\noindent \textit{दीर्घ उत्तरीय प्रश्न। लगभग 250 शब्दों में उत्तर दीजिए।}

\begin{enumerate}[label=\textbf{\arabic*.},leftmargin=*,itemsep=12pt]

  \item What is the Sangam Age? What is our political, social and economic understanding about this period from literary sources?\pts{10}
  
  संगम युग क्या है? साहित्यिक स्रोतों के आधार पर इस काल के विषय में हमारी राजनीतिक, सामाजिक तथा आर्थिक समझ क्या है?

  \begin{center}\textbf{OR / अथवा}\end{center}

  Discuss the political achievements and administrative development of the Imperial Chola rulers, especially Rajaraja I and Rajendra I.\pts{10}
  
  चोल शासकों, विशेष रूप से राजराज प्रथम तथा राजेन्द्र प्रथम की राजनीतिक उपलब्धियों एवं प्रशासनिक विकास की चर्चा कीजिए।

  \item Describe the political growth and cultural contributions of the Pandyan rulers, with special reference to the early and later Pandyas.\pts{10}
  
  प्रारंभिक तथा उत्तरकालीन पाण्ड्य शासकों के विशेष संदर्भ में पाण्ड्य राजाओं के राजनीतिक विकास एवं सांस्कृतिक योगदान का वर्णन कीजिए।

  \begin{center}\textbf{OR / अथवा}\end{center}

  Compare the political and cultural contributions of the Satavahana and Vakataka dynasties, highlighting their role in the history of the Deccan.\pts{10}
  
  दक्कन के इतिहास में सातवाहन तथा वाकाटक वंशों की भूमिका को रेखांकित करते हुए उनके राजनीतिक एवं सांस्कृतिक योगदान की तुलना कीजिए।

  \item Examine the major rulers of the Pallava dynasty and evaluate their political and cultural contributions that established their distinct place in the history of Southern India.\pts{10}
  
  पल्लव वंश के प्रमुख शासकों का परीक्षण कीजिए तथा उनके राजनीतिक एवं सांस्कृतिक योगदान का मूल्यांकन कीजिए।

\end{enumerate}

\bigskip

\section*{SECTION -- B / खण्ड -- ब}
\noindent \textit{Note: Short Answer Type Questions. Answer approximately in 125--150 words.}\pts{4 $\times$ 5 = 20}

\noindent \textit{लघु उत्तरीय प्रश्न। लगभग 125--150 शब्दों में उत्तर दीजिए।}

\begin{enumerate}[label=\textbf{\arabic*.},leftmargin=*,start=4,itemsep=10pt]

  \item Write a short paragraph on Tailapa II and explain how he revived Chalukya power in the Deccan.\pts{5}
  
  तैलप द्वितीय पर एक संक्षिप्त टिप्पणी लिखिए तथा व्याख्या कीजिए कि उसने दक्कन में चालुक्य शक्ति का पुनरुद्धार कैसे किया।

  \begin{center}\textbf{OR / अथवा}\end{center}

  Briefly explain how Vikramaditya I revived Chalukya power after the Pallava invasions.\pts{5}
  
  पल्लव आक्रमणों के पश्चात् विक्रमादित्य प्रथम ने चालुक्य सत्ता का पुनरुद्धार किस प्रकार किया, संक्षेप में समझाइए।

  \item Write a short paragraph on Vakataka--Gupta relations and their historical importance.\pts{5}
  
  वाकाटक-गुप्त सम्बन्धों तथा उनके ऐतिहासिक महत्त्व पर एक संक्षिप्त टिप्पणी लिखिए।

  \begin{center}\textbf{OR / अथवा}\end{center}

  Explain the importance of inscriptions and coins in understanding Satavahana history.\pts{5}
  
  सातवाहन इतिहास को समझने में अभिलेखों तथा सिक्कों के महत्त्व की व्याख्या कीजिए।

  \item Briefly describe the political achievements of Gautamiputra Satakarni.\pts{5}
  
  गौतमीपुत्र शातकर्णी की राजनीतिक उपलब्धियों का संक्षेप में वर्णन कीजिए।

  \begin{center}\textbf{OR / अथवा}\end{center}

  Write a short note on Dantidurga and the rise of the Rashtrakuta dynasty.\pts{5}
  
  दन्तिदुर्ग तथा राष्ट्रकूट वंश के अभ्युदय पर एक संक्षिप्त टिप्पणी लिखिए।

\end{enumerate}

\bigskip

\section*{SECTION -- C / खण्ड -- स}
\noindent \textit{Note: Very Short Answer Type Questions. Answer approximately in 15--20 words.}\pts{10 $\times$ 2 = 20}

\noindent \textit{अति लघु उत्तरीय प्रश्न। लगभग 15--20 शब्दों में उत्तर दीजिए।}

\begin{enumerate}[label=\textbf{\arabic*.},leftmargin=*,start=7,itemsep=8pt]
  \item 
  \begin{parts}
    \item Which Vakataka queen was the daughter of Chandragupta II?\pts{2}
    
    कौन-सी वाकाटक रानी चंद्रगुप्त द्वितीय की पुत्री थी? (Prabhavatigupta / प्रभावतीगुप्त)

    \item Which Chalukya king defeated King Harsha?\pts{2}
    
    किस चालुक्य राजा ने राजा हर्ष को पराजित किया था? (Pulakeshin II / पुलकेशिन द्वितीय)

    \item Which Rashtrakuta king built the Kailasa temple at Ellora?\pts{2}
    
    एलोरा के कैलाश मन्दिर का निर्माण किस राष्ट्रकूट राजा ने करवाया था? (Krishna I / कृष्ण प्रथम)

    \item What was the capital of the Yadava dynasty?\pts{2}
    
    यादव वंश की राजधानी क्या थी? (Devagiri / देवगिरि)

    \item Which Rashtrakuta ruler composed the Kannada work `Kavirajamarga'?\pts{2}
    
    कन्नड़ ग्रन्थ `कविराजमार्ग' की रचना किस राष्ट्रकूट शासक ने की थी? (Amoghavarsha I / अमोघवर्ष प्रथम)

    \item Which Vakataka king held the title ``Samrat'' and expanded the empire widely?\pts{2}
    
    किस वाकाटक राजा ने `सम्राट' की उपाधि धारण की थी? (Pravarasena I / प्रवरसेन प्रथम)

    \item Which Satavahana king issued coins with ship symbols showing maritime trade?\pts{2}
    
    किस सातवाहन राजा ने समुद्री व्यापार दर्शाने वाले जहाजी सिक्के जारी किए थे? (Yajna Sri Satakarni / यज्ञ श्री शातकर्णी)

    \item Which Pallava ruler built the rock-cut temples (Ratha temples) at Mahabalipuram?\pts{2}
    
    महाबलीपुरम के रथ मन्दिरों का निर्माण किस पल्लव शासक ने करवाया था? (Narasimhavarman I / नरसिंहवर्मन प्रथम)

    \item Who founded the Hoysala dynasty?\pts{2}
    
    होयसल वंश का संस्थापक कौन था? (Nripa Kama II / Sala)

    \item Name the naval capital of the Cholas under Rajendra I.\pts{2}
    
    राजेन्द्र प्रथम के शासनकाल में चोलों की नौसैनिक राजधानी/प्रसिद्ध बन्दरगाह का नाम लिखिये। (Nagapattinam / Gangaikonda Cholapuram)
  \end{parts}
\end{enumerate}

\end{document}
